\documentclass[12pt, a4paper]{article}

% ── PACKAGES ──────────────────────────────────────────────────
\usepackage[utf8]{inputenc}
\usepackage[T1]{fontenc}
\usepackage{geometry}
\geometry{margin=2.5cm}
\usepackage{hyperref}
\hypersetup{
    colorlinks=true,
    linkcolor=blue,
    urlcolor=blue,
    citecolor=blue,
    pdftitle={Tazemetostat Followed by Fulvestrant in
              EZH2-High FOXA1-Absent TNBC},
    pdfauthor={Eric Robert Lawson},
    pdfkeywords={tazemetostat, fulvestrant, TNBC,
                 EZH2, FOXA1, PRC2, breast cancer,
                 epigenetic reprogramming, trial proposal,
                 OrganismCore, attractor geometry}
}
\usepackage{amsmath}
\usepackage{booktabs}
\usepackage{longtable}
\usepackage{array}
\usepackage{parskip}
\usepackage{microtype}
\usepackage{authblk}
\usepackage{titlesec}
\usepackage{fancyhdr}
\usepackage{xcolor}
\usepackage{float}
\usepackage{enumitem}
\usepackage{setspace}

% ── COLOURS ───────────────────────────────────────────────────
\definecolor{repobox}{RGB}{240,247,255}
\definecolor{urgentbox}{RGB}{255,235,235}
\definecolor{lockbox}{RGB}{245,255,245}
\definecolor{mechanismbox}{RGB}{250,245,255}

% ── HEADER / FOOTER ───────────────────────────────────────────
\pagestyle{fancy}
\fancyhf{}
\rhead{\small OrganismCore \textbar{} 2026-03-06}
\lhead{\small CS-LIT-16: Tazemetostat $\rightarrow$
       Fulvestrant in TNBC}
\rfoot{\small Page \thepage}
\renewcommand{\headrulewidth}{0.4pt}

% ── SECTION FORMATTING ────────────────────────────────────────
\titleformat{\section}
  {\large\bfseries}{\thesection}{1em}{}[\titlerule]
\titleformat{\subsection}
  {\normalsize\bfseries}{\thesubsection}{1em}{}

% ══════════════════════════════════════════════════════════════
% TITLE
% ══════════════════════════════════════════════════════════════
\title{
    \vspace{-1em}
    \textbf{Tazemetostat Followed by Fulvestrant in
    EZH2-High, FOXA1-Absent Triple-Negative Breast
    Cancer:}\\[0.5em]
    \large A Principles-First Trial Proposal with
    Mechanistic Support from Two Independent
    Published Studies\\[0.4em]
    \normalsize \textit{OrganismCore Framework ---
    Document CS-LIT-16}\\[0.3em]
    \small Derived from Waddington Landscape Attractor
    Geometry Before Literature Review\\[0.3em]
    \small Confirmed Independently by Toska et al.\
    \textit{Nature Medicine} 2017 and
    Schade et al.\ \textit{Nature} 2024
}

\author[1]{Eric Robert Lawson}
\affil[1]{
    OrganismCore\\
    \href{mailto:OrganismCore@proton.me}
         {OrganismCore@proton.me}\\
    ORCID:
    \href{https://orcid.org/0009-0002-0414-6544}
         {0009-0002-0414-6544}
}

\date{March 6, 2026}

% ══════════════════════════════════════════════════════════════
\begin{document}
% ══════════════════════════════════════════════════════════════

\maketitle
\thispagestyle{fancy}

% ── URGENCY NOTICE ────────────────────────────────────────────
\begin{center}
\colorbox{urgentbox}{
\begin{minipage}{0.88\textwidth}
\vspace{0.5em}
\centering
\textbf{PRIORITY ESTABLISHMENT DOCUMENT}\\[0.3em]
\small
As of 2026-03-06, no registered clinical trial
exists anywhere in the world for tazemetostat
followed by fulvestrant in TNBC
(ClinicalTrials.gov searched 2026-03-05).
This document establishes priority on the specific
trial design, patient selection criteria, treatment
sequence, and primary endpoint described herein.
The mechanism was derived from first principles
before literature review and is confirmed by two
independent published papers cited below.
\vspace{0.5em}
\end{minipage}
}
\end{center}

\vspace{0.6em}

% ── REPOSITORY POINTER ────────────────────────────────────────
\begin{center}
\colorbox{repobox}{
\begin{minipage}{0.88\textwidth}
\vspace{0.4em}
\centering
\textbf{Full computational record and companion
reasoning artifacts:}\\[0.3em]
\href{https://github.com/Eric-Robert-Lawson/attractor-oncology}
{\texttt{https://github.com/Eric-Robert-Lawson/attractor-oncology}}
\\[0.2em]
\small
Repository ID: 1172048282 \textbar{}
MIT License \textbar{}
Python \textbar{}
Public\\[0.2em]
Companion deposit:
\textit{FOXA1/EZH2 Ratio as a Universal Breast Cancer
Subtype Classifier and Treatment Predictor}
(Zenodo, 2026-03-06)\\[0.2em]
Source document: \texttt{Cancer\_Research/BRCA/DEEP\_DIVE/}
\texttt{BRCA\_Cross\_Subtype\_Literature\_Check.md}
(BRCA-S8h, item CS-LIT-16)
\vspace{0.4em}
\end{minipage}
}
\end{center}

\vspace{0.8em}

% ── ABSTRACT ──────────────────────────────────────────────────
\begin{abstract}
\noindent
Triple-negative breast cancer (TNBC) accounts for
approximately 170,000 new diagnoses per year globally
and carries the worst prognosis of all breast cancer
subtypes. The apparent absence of oestrogen receptor
(ER) expression in TNBC has historically precluded
endocrine therapy. We propose that ER negativity in
EZH2-high TNBC is not a fixed biological state but
an epigenetically reversible condition: the luminal
identity programme --- including FOXA1, GATA3, and
ESR1 --- is silenced by PRC2/EZH2-mediated H3K27me3
deposition rather than being absent. This silencing
is pharmacologically reversible with tazemetostat
(EZH2 inhibitor, FDA-approved 2020). When EZH2
inhibition removes the silencing mark, FOXA1
re-expresses, luminal reprogramming initiates, and
ESR1 returns --- at which point fulvestrant
(selective ER degrader, FDA-approved 2002) can engage
the restored receptor.

This mechanistic chain was derived from Waddington
landscape attractor geometry applied to 19,542 single
cancer cells (GSE176078) before literature review.
It is independently confirmed by two published
studies: Toska et al.\ (\textit{Nature Medicine},
2017) demonstrated EZH2 inhibition re-expresses FOXA1
and initiates luminal reprogramming in ER-negative
breast cancer; Schade et al.\ (\textit{Nature}, 2024)
demonstrated PRC2/EZH2 inhibition drives basal-like
TNBC cells into a more differentiated, luminal-like
state with induction of luminal markers including
GATA3. Neither paper proposed the clinical sequence
described here.

Supporting computational validation: in GSE25066
($n=508$ neoadjuvant chemotherapy patients), the
EZH2 paradox is confirmed --- EZH2-high tumours are
chemotherapy-sensitive short-term (pCR, $p<0.0001$)
but carry worse distant relapse-free survival
long-term (HR$=1.363$, $p=0.0047$), identifying the
post-chemotherapy window as the critical intervention
point. The FOXA1/EZH2 ratio correctly classifies TNBC
from luminal subtypes with AUC$=0.828$--$0.901$
across two independent datasets, providing the IHC
patient selection instrument.

We propose a Phase~I/II clinical trial of tazemetostat
($800$~mg BID $\times$ 4~weeks) followed by fulvestrant
($500$~mg q28d) in patients with EZH2-high
(H-score~$\geq150$), FOXA1-absent (H-score~$<20$)
TNBC. The primary endpoint is FOXA1 H-score
re-expression $\geq50$ at week~4 on-treatment biopsy.
No competing trial exists. Both drugs are FDA-approved.
The trial can be designed and submitted immediately.
\end{abstract}

\vspace{0.5em}
\noindent\textbf{Keywords:}
tazemetostat; fulvestrant; TNBC; triple-negative
breast cancer; EZH2; FOXA1; PRC2; epigenetic
reprogramming; luminal reprogramming; H3K27me3;
ESR1; GATA3; endocrine therapy; EZH2 inhibitor;
attractor geometry; Waddington landscape; OrganismCore;
trial proposal; patient selection; IHC

\newpage
\tableofcontents
\newpage

% ══════════════════════════════════════════════════════════════
\section{Document Metadata}
% ══════════════════════════════════════════════════════════════

\begin{center}
\begin{tabular}{ll}
\toprule
\textbf{Field} & \textbf{Value} \\
\midrule
Document ID       & CS-LIT-16-TRIAL-PROPOSAL \\
Series            & BRCA Deep Dive --- Cross-Subtype \\
Type              & Clinical Trial Proposal ---
                    Priority Establishment Document \\
Parent document   & BRCA-S8h (Cross-Subtype Literature
                    Check) \\
Companion deposit & FOXA1/EZH2 Ratio Master Reasoning
                    Artifact (Zenodo, 2026-03-06) \\
Date              & 2026-03-06 \\
Author            & Eric Robert Lawson \\
Organisation      & OrganismCore \\
ORCID             & 0009-0002-0414-6544 \\
Status            & Permanent \\
Repository        &
  \href{https://github.com/Eric-Robert-Lawson/attractor-oncology}
       {Eric-Robert-Lawson/attractor-oncology} \\
Repository ID     & 1172048282 \\
License           & CC BY 4.0 \\
Biological contradictions & 0 \\
Registered trials & 0 (as of 2026-03-05) \\
\bottomrule
\end{tabular}
\end{center}

% ══════════════════════════════════════════════════════════════
\section{Background: The Problem with TNBC}
% ══════════════════════════════════════════════════════════════

\subsection{Current treatment landscape and its limits}

Triple-negative breast cancer (TNBC) is defined by the
absence of oestrogen receptor (ER), progesterone receptor
(PR), and HER2 amplification. It accounts for approximately
15\% of all breast cancers --- roughly 170,000 new diagnoses
per year globally --- and disproportionately affects younger
women and women of African ancestry. Five-year survival for
metastatic TNBC remains below 30\%.

Current standard-of-care for TNBC is:
\begin{itemize}[noitemsep]
  \item Neoadjuvant taxane-anthracycline chemotherapy
  \item Pembrolizumab (PD-1 checkpoint inhibitor) in
        PD-L1+ or high-TIL cases
  \item Olaparib/talazoparib in BRCA1/2-mutant cases
  \item Sacituzumab govitecan in metastatic setting
\end{itemize}

Endocrine therapy is not used in TNBC. The universal
assumption is that ER-negative tumours cannot respond to
endocrine agents because there is no receptor to target.
This assumption has never been challenged with an EZH2
inhibitor priming approach.

\subsection{Why the assumption may be wrong for EZH2-high TNBC}

ER negativity in TNBC is not uniform in its molecular basis.
Two mechanistically distinct causes of ER negativity exist:

\begin{center}
\begin{tabular}{p{3cm}p{5.5cm}p{5cm}}
\toprule
\textbf{Type} & \textbf{Mechanism} & \textbf{Implication} \\
\midrule
Commitment-absent
(Claudin-low)
& The luminal identity programme was
  never activated. FOXA1, GATA3,
  ESR1 are absent because the cell
  never committed to luminal lineage.
& No ER to restore.
  Endocrine therapy cannot work.
  Immune compartment is the only target. \\[0.5em]
Epigenetically
silenced
(EZH2-high TNBC)
& The luminal identity programme is
  present but silenced by PRC2/EZH2
  H3K27me3 marks on FOXA1, GATA3,
  ESR1 promoters.
& ER can potentially be restored.
  If EZH2 inhibition removes the
  silencing mark, the luminal
  programme re-activates and
  endocrine therapy can engage. \\
\bottomrule
\end{tabular}
\end{center}

The EZH2 paradox --- quantified in this framework in
GSE25066 (see Section~\ref{sec:ezh2paradox}) --- confirms
that EZH2-high TNBC represents a distinct biological
state where the luminal programme is suppressed rather
than absent.

% ══════════════════════════════════════════════════════════════
\section{The Mechanistic Case}
% ══════════════════════════════════════════════════════════════

\subsection{The Waddington geometry derivation}

The OrganismCore framework applies Waddington landscape
attractor geometry to single-cell RNA-seq data from real
tumour biopsies. In this model, cancer subtypes correspond
to attractor states --- stable configurations of gene
expression maintained by specific molecular locks.

Applied to 19,542 single cancer cells from 26 patients
(GSE176078, Wu et al.\ \textit{Nature Genetics} 2021),
the framework identified TNBC as an attractor state
characterised by:

\begin{itemize}[noitemsep]
  \item EZH2 elevated $+189\%$ above normal Mature
        Luminal reference
  \item FOXA1 at near-zero expression (epigenetically
        silenced, not absent from the genome)
  \item ESR1 silenced downstream of FOXA1 suppression
  \item GATA3 silenced by the same H3K27me3 mechanism
\end{itemize}

The geometry identified the lock as PRC2/EZH2-mediated
H3K27me3 deposition. The predicted unlock sequence was:

\begin{center}
\colorbox{mechanismbox}{
\begin{minipage}{0.82\textwidth}
\vspace{0.5em}
\centering
\textbf{Step 1:} Tazemetostat (EZH2 inhibitor)\\
$\downarrow$\\
EZH2 catalytic activity inhibited\\
H3K27me3 marks removed from FOXA1/GATA3/ESR1 promoters\\
$\downarrow$\\
\textbf{Step 2:} FOXA1 re-expresses\\
$\downarrow$\\
FOXA1 pioneer factor opens luminal chromatin\\
GATA3 and ESR1 transcription restored\\
$\downarrow$\\
\textbf{Step 3:} Fulvestrant (selective ER degrader)\\
$\downarrow$\\
Engages the restored oestrogen receptor\\
Degrades ER protein, removing the luminal growth signal\\
$\downarrow$\\
\textbf{Outcome:} TNBC converted to endocrine-responsive
state and treated accordingly
\vspace{0.5em}
\end{minipage}
}
\end{center}

This sequence was derived from geometry before any
literature review. The literature review confirmed it
independently from two directions.

\subsection{Independent published confirmation}

\begin{center}
\begin{tabular}{p{3.5cm}p{9.5cm}}
\toprule
\textbf{Paper} & \textbf{Finding relevant to this proposal} \\
\midrule
Toska et al.\\
\textit{Nature Medicine}\\
2017
&
EZH2 inhibition in ER-negative breast cancer
re-expresses FOXA1. FOXA1 acts as a pioneer factor
opening luminal chromatin. Cells acquire a luminal-like
transcriptional programme. The authors conclude that
EZH2 inhibitors can promote FOXA1 re-expression and
luminal reprogramming in ER-negative breast cancer.
\textbf{This is Step 1--2 of the proposed sequence,
independently confirmed.}
\\[0.5em]
Schade et al.\\
\textit{Nature}\\
2024\\
(Harvard/Ludwig)
&
PRC2/EZH2 inhibition drives basal-like TNBC cells into
a more differentiated, luminal-like state. Molecular
profiling confirms induction of luminal markers
including GATA3 after EZH2 inhibition. The paper
identifies EZH2/PRC2 as the convergence node in
basal-like TNBC that silences FOXA1/GATA3/ESR1.
\textbf{This is independent confirmation of the
entire mechanistic premise.}
Schade et al.\ propose EZH2i followed by AKT inhibition
as one downstream strategy. The present proposal uses
EZH2i followed by fulvestrant as an alternative
downstream strategy targeting the restored ER.
Both are complementary, not competing.
\\
\bottomrule
\end{tabular}
\end{center}

\subsection{Why tazemetostat must precede fulvestrant}

This sequencing is not interchangeable. The order is
mechanistically required:

\begin{enumerate}
  \item Fulvestrant degrades the oestrogen receptor
        protein. It requires an expressed ER to act on.
        In EZH2-high TNBC, ESR1 is epigenetically
        silenced. There is no expressed ER protein.
        Fulvestrant given first has no target.
        It will show no activity in this population ---
        which is consistent with all published data
        showing fulvestrant is inactive in ER-negative
        breast cancer.

  \item Tazemetostat removes the H3K27me3 silencing
        marks from ESR1 and FOXA1. After 3--6 weeks
        of EZH2 inhibition, FOXA1 re-expresses.
        FOXA1 opens chromatin. ESR1 transcription
        is restored. ER protein is produced.

  \item Once ER protein is present, fulvestrant can
        engage it, bind, and degrade it --- removing
        the luminal growth signal.

  \item The critical biomarker bridge between steps:
        FOXA1 H-score on on-treatment biopsy at week~4
        confirms the conversion has occurred before
        fulvestrant is added. This is the primary
        endpoint of the proposed trial.
\end{enumerate}

% ══════════════════════════════════════════════════════════════
\section{Supporting Computational Evidence}
\label{sec:ezh2paradox}
% ══════════════════════════════════════════════════════════════

\subsection{The EZH2 paradox: two windows, same biology}

The EZH2 paradox was identified and quantified in
GSE25066 (Hatzis et al.\ \textit{JAMA} 2011, $n=508$
neoadjuvant chemotherapy patients). It has two arms
that appear contradictory but are mechanistically unified:

\begin{center}
\begin{tabular}{p{3.5cm}p{4.5cm}p{4.5cm}}
\toprule
& \textbf{Short-window arm} & \textbf{Long-window arm} \\
\midrule
\textbf{Finding}
& EZH2-high tumours achieve
  pathological complete response
  (pCR) more often with
  neoadjuvant chemotherapy
& EZH2-high tumours have worse
  distant relapse-free survival
  after 3--5 years \\[0.3em]
\textbf{Statistic}
& pCR $p < 0.0001$
  (GSE25066, $n=508$)
& HR $= 1.363$, $p = 0.0047$
  (GSE25066 DRFS) \\[0.3em]
\textbf{Mechanism}
& EZH2 drives proliferation.
  High proliferation $=$
  chemo-sensitive.
  Chemotherapy kills dividing cells.
& EZH2 silences tumour suppression.
  Residual EZH2-high cells after
  chemo escape as late-relapsing
  metastatic disease. \\[0.3em]
\textbf{Published support}
& Logically inferrable from
  published biology
& Frontiers Oncology 2020
  ($\sim$4,000 patients);
  Fineberg (Montefiore):
  EZH2 IHC predicts metastatic
  TNBC post-neoadjuvant chemo \\
\bottomrule
\end{tabular}
\end{center}

\noindent\textbf{What the paradox reveals:}
EZH2-high TNBC patients often achieve pCR with
chemotherapy. But pCR is not cure in this population.
The EZH2-high cells that survive chemotherapy ---
either residual disease or circulating tumour cells
--- carry the same epigenetic programme that silences
luminal identity. They are the source of late relapse.

The intervention window is the period between
chemotherapy completion and late relapse
(approximately 3--5 years). The tazemetostat
$\rightarrow$ fulvestrant sequence is designed for
this window: in patients who achieve pCR or have
minimal residual disease after chemotherapy, converting
the remaining EZH2-high TNBC cell population to an
endocrine-responsive state and treating it accordingly.

\subsection{FOXA1/EZH2 ratio validation data}

The FOXA1/EZH2 ratio, validated across $\sim$7,500
patients in seven independent datasets (companion
Zenodo deposit, 2026-03-06), provides the patient
selection instrument for this trial. Key figures:

\begin{center}
\begin{tabular}{lcc}
\toprule
\textbf{Classification task} & \textbf{Dataset}
& \textbf{AUC} \\
\midrule
LumA vs Basal/TNBC & METABRIC ($n=909$) & 0.828 \\
LumA vs Basal/TNBC & TCGA PanCan ($n=151$) & 0.901 \\
\bottomrule
\end{tabular}
\end{center}

The ratio correctly identifies TNBC (low ratio) from
luminal subtypes (high ratio) with AUC$=0.83$--$0.90$
replicated in two independent datasets. This AUC range,
combined with EZH2-specific IHC selection criteria
(H-score $\geq 150$), provides a robust patient
identification instrument for trial enrolment.

\subsection{EZH2 graded elevation: the biological basis
for patient selection}

From single-cell data (GSE176078, $n=19{,}542$ cancer
cells), EZH2 expression in TNBC is elevated $+189\%$
above normal Mature Luminal reference. This is confirmed
in published literature:

\begin{itemize}[noitemsep]
  \item Frontiers in Oncology 2020: EZH2/NSD2 axis
        predicts poor prognosis in $\sim$4,000 breast
        cancer patients; TNBC highest.
  \item Schade et al.\ \textit{Nature} 2024: EZH2/PRC2
        is the explicitly named convergence node in
        basal-like TNBC (CONFIRMED, CS-LIT-4).
\end{itemize}

Not all TNBC is equally EZH2-high. The patient selection
criterion (H-score $\geq 150$) enriches for the
subpopulation with the highest probability of PRC2-driven
FOXA1 silencing and therefore the highest probability
of FOXA1 re-expression upon tazemetostat treatment.

% ══════════════════════════════════════════════════════════════
\section{The Clinical Trial Proposal}
% ══════════════════════════════════════════════════════════════

\subsection{Trial overview}

\begin{center}
\begin{tabular}{ll}
\toprule
\textbf{Parameter} & \textbf{Specification} \\
\midrule
Phase & I/II, open-label, single-arm \\
Setting & Metastatic or locally advanced TNBC;
          also post-neoadjuvant window \\
Design & Sequential: tazemetostat alone
         $\times$ 4 weeks, then add fulvestrant \\
Primary endpoint & FOXA1 H-score re-expression
                   $\geq 50$ at week 4 biopsy \\
Sample size (Phase I) & $n = 12$--$18$
                        (safety + biomarker) \\
Sample size (Phase II) & $n = 40$--$60$
                         (efficacy signal) \\
Estimated duration & 18--24 months \\
Infrastructure & Any oncology centre with
                 IHC capability \\
Both drugs approved & Yes (FDA-approved) \\
\bottomrule
\end{tabular}
\end{center}

\subsection{Patient selection criteria}

\subsubsection{Inclusion criteria}

\begin{enumerate}
  \item Histologically confirmed TNBC
        (ER $<1\%$, PR $<1\%$, HER2 IHC $0$--$1+$
        or FISH non-amplified)
  \item EZH2 H-score $\geq 150$ on diagnostic or
        recent biopsy (EZH2 IHC, validated antibody)
  \item FOXA1 H-score $< 20$ on diagnostic or
        recent biopsy (FOXA1 IHC, validated antibody)
  \item Measurable disease (RECIST 1.1) or
        minimal residual disease post-neoadjuvant
        chemotherapy
  \item ECOG performance status $0$--$2$
  \item Adequate organ function
  \item Prior chemotherapy permitted
        (taxane-anthracycline standard of care)
  \item Willingness to undergo on-treatment biopsy
        at week 4
\end{enumerate}

\subsubsection{Exclusion criteria}

\begin{enumerate}
  \item FOXA1 H-score $\geq 20$ at baseline
        (conversion already present or non-EZH2-driven
        subtype)
  \item Prior EZH2 inhibitor therapy
  \item Prior fulvestrant therapy
  \item BRCA1/2 germline mutation with available
        PARP inhibitor (standard treatment available)
  \item Active CNS metastases requiring treatment
  \item Concurrent strong CYP3A4 inhibitors
        (tazemetostat interaction)
\end{enumerate}

\subsection{Treatment protocol}

\begin{center}
\colorbox{lockbox}{
\begin{minipage}{0.85\textwidth}
\vspace{0.5em}
\textbf{TREATMENT SEQUENCE (mechanistically required order)}

\medskip
\textbf{Phase 1 --- EZH2 Priming (Weeks 1--4):}\\
Tazemetostat 800~mg orally twice daily (BID)\\
Standard approved dose (follicular lymphoma/epithelioid
sarcoma)\\
Duration: 4 weeks continuous

\medskip
\textbf{On-treatment biopsy at Week 4:}\\
FOXA1 H-score measured\\
If FOXA1 H-score $\geq 50$: proceed to Phase 2\\
If FOXA1 H-score $< 50$: continue tazemetostat
$\times$ 4 additional weeks,
repeat biopsy at Week 8

\medskip
\textbf{Phase 2 --- Endocrine Engagement:}\\
Tazemetostat 800~mg BID (continue)\\
\textit{plus}\\
Fulvestrant 500~mg intramuscular injection on Days 1,
15, 29, then q28d\\
Standard approved dose\\
Continue until progression or unacceptable toxicity

\medskip
\textbf{Rationale for continuation of tazemetostat:}\\
EZH2 inhibition must be maintained while fulvestrant
acts. If tazemetostat is stopped, H3K27me3 marks
will re-accumulate, FOXA1 will be re-silenced,
and ER expression will be lost again.
Continuous EZH2 inhibition is required to maintain
the converted state.
\vspace{0.5em}
\end{minipage}
}
\end{center}

\subsection{Endpoints}

\subsubsection{Primary endpoint}

\begin{itemize}
  \item \textbf{FOXA1 H-score re-expression $\geq 50$
        at Week~4 on-treatment biopsy}
  \item Definition of response: any FOXA1 H-score
        $\geq 50$ in a patient with baseline H-score
        $< 20$
  \item Measured by: FOXA1 IHC (validated antibody
        clone, same laboratory as baseline)
  \item This is a biological activity endpoint, not
        a tumour shrinkage endpoint
  \item Justification: FOXA1 re-expression is the
        necessary intermediate step between EZH2
        inhibition and ER restoration. If FOXA1
        does not re-express, fulvestrant has no
        biological rationale to add. The trial
        tests whether the mechanistic chain
        (EZH2i $\rightarrow$ FOXA1) is operable
        in clinical tissue before proceeding to
        the downstream step.
\end{itemize}

\subsubsection{Key secondary endpoints}

\begin{enumerate}
  \item ESR1 re-expression at Week 4 and Week 8
        (ER IHC H-score)
  \item GATA3 re-expression at Week 4 and Week 8
  \item Objective response rate (ORR) per RECIST 1.1
        at Week 12 and Week 24
  \item Progression-free survival (PFS)
  \item Overall survival (OS)
  \item Distant relapse-free survival (DRFS)
        at 3 years
  \item H3K27me3 reduction in on-treatment biopsy
        (pharmacodynamic confirmation of EZH2 target
        engagement)
  \item Safety and tolerability profile of the
        combination
\end{enumerate}

\subsubsection{Exploratory endpoints}

\begin{enumerate}
  \item Correlation between baseline EZH2 H-score
        and FOXA1 re-expression rate
  \item Correlation between FOXA1 H-score at Week 4
        and ORR at Week 12
  \item ctDNA ESR1 dynamics during treatment
  \item Transcriptomic profiling (serial biopsies):
        luminal gene programme re-expression trajectory
\end{enumerate}

\subsection{Biomarker strategy}

\begin{center}
\begin{tabular}{p{3cm}p{4cm}p{5.5cm}}
\toprule
\textbf{Timepoint} & \textbf{Tissue collection}
& \textbf{Assays} \\
\midrule
Baseline (screening)
& Core needle biopsy
& EZH2 IHC (H-score)\\
& & FOXA1 IHC (H-score)\\
& & ER IHC (H-score)\\
& & GATA3 IHC\\
& & H3K27me3 IHC\\
& & Ki67\\
& & RNA-seq (if adequate tissue) \\[0.3em]
Week 4 (on-treatment)
& Core needle biopsy
& Same panel as baseline\\
& & \textbf{Primary endpoint measured here} \\[0.3em]
Week 8 (optional)
& Core needle biopsy
& Same panel (if Week 4 FOXA1 $<50$) \\[0.3em]
Week 12 (response)
& Imaging (RECIST)
& ctDNA (blood) \\[0.3em]
Progression
& Core needle biopsy (if accessible)
& Full panel + resistance mechanism profiling \\
\bottomrule
\end{tabular}
\end{center}

% ══════════════════════════════════════════════════════════════
\section{The Clinical Gap: Why This Trial Does Not Exist}
% ══════════════════════════════════════════════════════════════

\subsection{ClinicalTrials.gov search results}

Searched 2026-03-05. Search terms: ``tazemetostat
breast cancer'', ``tazemetostat TNBC'', ``tazemetostat
fulvestrant'', ``EZH2 inhibitor fulvestrant'',
``EZH2 inhibitor TNBC endocrine''.

\begin{center}
\begin{tabular}{ll}
\toprule
\textbf{Search} & \textbf{Result} \\
\midrule
Tazemetostat + breast cancer & No trials in TNBC \\
Tazemetostat + fulvestrant & No trials \\
Tazemetostat + TNBC & No trials \\
EZH2 inhibitor + fulvestrant & No trials \\
EZH2 inhibitor + TNBC & Preclinical only;
                         no clinical trials \\
\midrule
\textbf{Tazemetostat approved indications}
& EZH2-mutant follicular lymphoma (2020) \\
& Epithelioid sarcoma (2020) \\
\textbf{Active tazemetostat trials}
& ARID1A-mutant solid tumours
  (NCT05023655) \\
& NCI Pediatric MATCH basket \\
& None in breast cancer \\
\bottomrule
\end{tabular}
\end{center}

\noindent The clinical space for tazemetostat in TNBC
is completely unoccupied as of 2026-03-05.

\subsection{Why the gap exists}

The gap exists because the mechanistic logic requires
assembling three independent published findings into
a clinical sequence:

\begin{enumerate}
  \item EZH2 silences luminal identity in TNBC
        (Schade 2024 --- this is now published)
  \item EZH2 inhibition restores FOXA1 and luminal
        reprogramming in ER-negative breast cancer
        (Toska 2017 --- this is published)
  \item Therefore: use EZH2 inhibition to convert
        ER-negative TNBC to ER-positive, then apply
        endocrine therapy to the converted state
        (this assembly is not published anywhere)
\end{enumerate}

Step 3 does not follow automatically from steps 1
and 2 in the published literature. No paper makes
this connection. The assembly requires recognising
that the sequence is not just biologically possible
but clinically actionable using two approved drugs
with known safety profiles. That recognition is the
contribution of the OrganismCore attractor geometry
framework. It identified the clinical sequence from
the geometry of the attractor states before the
individual papers were reviewed.

\subsection{The Schade 2024 comparison}

Schade et al.\ (\textit{Nature} 2024) is the most
directly relevant published paper. It independently
confirms the central premise (EZH2i drives TNBC toward
luminal differentiation). Their proposed downstream
strategy differs from ours:

\begin{center}
\begin{tabular}{p{3cm}p{5.5cm}p{5cm}}
\toprule
& \textbf{Schade et al.\ 2024} & \textbf{This proposal} \\
\midrule
Step 1 (shared)
& EZH2i drives luminal differentiation
& EZH2i drives luminal differentiation \\
Step 2 (diverges)
& AKT inhibition to exploit the
  differentiated state's new
  dependencies (involution-like
  cell death)
& Fulvestrant to engage the restored
  ER and degrade it (endocrine therapy
  on the converted state) \\
Mechanism of kill
& Synthetic lethality with AKT
  in the differentiated state
& ER degradation in the converted
  luminal state \\
Drug status
& AKT inhibitor (capivasertib:
  FDA-approved for HR+ HER2-low;
  not approved in TNBC)
& Fulvestrant (FDA-approved;
  standard formulary globally) \\
\bottomrule
\end{tabular}
\end{center}

These two strategies are complementary. They can be
tested in parallel. Ours is distinguished by using
fully approved, universally available, and
globally affordable drugs in the second step.

% ══════════════════════════════════════════════════════════════
\section{Derivation Provenance}
% ══════════════════════════════════════════════════════════════

\subsection{Before-document confirmation}

The OrganismCore framework operates under a strict
before-document protocol: all predictions are locked
in version-controlled documents before any confirmatory
analysis or literature review is performed. The
tazemetostat $\rightarrow$ fulvestrant sequence for
TNBC was derived from attractor geometry and recorded
in the framework's reasoning artifacts before the
literature check (BRCA-S8h) was conducted.

The derivation sequence in the repository:

\begin{center}
\begin{tabular}{lll}
\toprule
\textbf{Document} & \textbf{Content} & \textbf{Status} \\
\midrule
BRCA-S8c-PLAIN & Script 1 plain account
                & Geometry derivation \\
BRCA-S8e-PLAIN & Script 2 update
                & Ratio confirmed \\
BRCA-S8g       & Script 3 results
                & Lock types documented \\
BRCA-S8h       & Literature check
                & CS-LIT-16 confirmed as
                  CONVERGENT-NOVEL \\
CS-LIT-16      & This document
                & Trial proposal locked \\
\bottomrule
\end{tabular}
\end{center}

All documents are version-controlled with cryptographic
timestamps in the repository:
\href{https://github.com/Eric-Robert-Lawson/attractor-oncology}
{\texttt{github.com/Eric-Robert-Lawson/attractor-oncology}}

\subsection{Literature check verdict: CONVERGENT-NOVEL}

From BRCA-S8h (Cross-Subtype Literature Check,
2026-03-05), item CS-LIT-16 received the verdict
\textbf{CONVERGENT-NOVEL}:

\begin{quote}
\itshape
``The mechanism (EZH2i $\rightarrow$ FOXA1
re-expression $\rightarrow$ luminal reprogramming)
is confirmed by two independent published papers
from two independent groups (Toska 2017, Schade 2024).
These papers did not know about each other's full work
when published and did not know about this framework.
They all found the same biological step.

\upshape\itshape
The clinical application --- using this mechanism as
the basis for a two-drug sequence (tazemetostat THEN
fulvestrant) in a specific patient population
(EZH2-high, FOXA1-absent TNBC) with a specific
biomarker endpoint (FOXA1 re-expression on biopsy
at week 4) --- has not been proposed or trialled.

The pieces are confirmed. The assembly is novel.
The trial does not exist. The clinical space is empty.

This is the most actionable CONVERGENT-NOVEL item
in the framework.''
\upshape
\end{quote}

% ══════════════════════════════════════════════════════════════
\section{The 30-Finding Context}
% ══════════════════════════════════════════════════════════════

This proposal is one of 30 findings from the BRCA
cross-subtype analysis (BRCA-S8h). It is classified
as the highest clinical priority for the following
reasons:

\begin{center}
\begin{tabular}{p{5cm}p{8cm}}
\toprule
\textbf{Criterion} & \textbf{Status for CS-LIT-16} \\
\midrule
Mechanism confirmed independently
& Yes --- two Nature-tier papers \\
Clinical space empty
& Yes --- no registered trial worldwide \\
Both drugs approved
& Yes --- FDA-approved, standard formulary \\
Patient selection instrument available
& Yes --- EZH2/FOXA1 IHC (standard stains) \\
Primary endpoint measurable
& Yes --- FOXA1 H-score on biopsy \\
Patient population defined
& Yes --- EZH2-high ($\geq150$),
  FOXA1-absent ($<20$) TNBC \\
Number of patients affected
& $\sim$170,000 new TNBC per year globally \\
Unmet medical need
& Severe --- 5-year metastatic survival $<30\%$ \\
\bottomrule
\end{tabular}
\end{center}

No other item in the 30-finding framework meets all
eight criteria simultaneously.

% ══════════════════════════════════════════════════════════════
\section{Locked Priority Statement}
% ══════════════════════════════════════════════════════════════

\begin{center}
\colorbox{urgentbox}{
\begin{minipage}{0.88\textwidth}
\vspace{0.6em}
\textbf{LOCKED AS OF 2026-03-06}

\medskip
\textbf{The proposal:}
Tazemetostat (EZH2 inhibitor, 800~mg BID $\times$
4~weeks) followed by fulvestrant (500~mg q28d) in
patients with EZH2-high (H-score $\geq 150$),
FOXA1-absent (H-score $< 20$) triple-negative
breast cancer.

\medskip
\textbf{The primary endpoint:}
FOXA1 H-score re-expression $\geq 50$ at Week~4
on-treatment biopsy.

\medskip
\textbf{The mechanistic basis:}
PRC2/EZH2 silences FOXA1, GATA3, and ESR1 in TNBC.
Tazemetostat removes H3K27me3 marks.
FOXA1 re-expresses. ESR1 is restored.
Fulvestrant engages and degrades the restored ER.

\medskip
\textbf{Independent confirmation:}
Toska et al.\ \textit{Nature Medicine} 2017 ---
EZH2i $\rightarrow$ FOXA1 re-expression in
ER-negative breast cancer.\\
Schade et al.\ \textit{Nature} 2024 ---
EZH2i drives TNBC toward luminal-like state.

\medskip
\textbf{The gap:}
No registered trial as of 2026-03-06.
Both drugs FDA-approved.
The trial can be designed and submitted today.

\medskip
\textbf{Derived from:}
OrganismCore Waddington landscape attractor geometry.
Prediction locked before literature review.
Repository: \href{https://github.com/Eric-Robert-Lawson/attractor-oncology}
{\texttt{github.com/Eric-Robert-Lawson/attractor-oncology}}

\medskip
\hfill --- Eric Robert Lawson, March 6, 2026
\vspace{0.4em}
\end{minipage}
}
\end{center}

% ══════════════════════════════════════════════════════════════
\section*{Repository and Reproducibility}
\addcontentsline{toc}{section}{Repository and
Reproducibility}
% ═��════════════════════════════════════════════════════════════

All reasoning artifacts, computational scripts, and
supporting data are available in the public repository:

\begin{center}
\href{https://github.com/Eric-Robert-Lawson/attractor-oncology}
{\texttt{https://github.com/Eric-Robert-Lawson/attractor-oncology}}\\[0.3em]
\small Repository ID: 1172048282 \textbar{}
MIT License \textbar{}
Python \textbar{}
Public
\end{center}

\noindent Relevant files:

\begin{center}
\begin{tabular}{ll}
\toprule
\textbf{File} & \textbf{Path} \\
\midrule
Literature check (parent document)
& \texttt{Cancer\_Research/BRCA/DEEP\_DIVE/}\\
& \texttt{BRCA\_Cross\_Subtype\_Literature\_Check.md} \\
FOXA1/EZH2 ratio master artifact
& \texttt{.../FOXA1\_EZH2\_RATIO/}\\
& \texttt{Master\_Artifact\_After\_Scripts.md} \\
Script 4 (EZH2 paradox, GSE25066)
& \texttt{.../FOXA1\_EZH2\_RATIO/ratio4.py} \\
Script 4 log (raw results)
& \texttt{.../ratio\_s4\_results/results/}\\
& \texttt{ratio\_s4\_log.txt} \\
\bottomrule
\end{tabular}
\end{center}

% ══════════════════════════════════════════════════════════════
\section*{Author Statement}
\addcontentsline{toc}{section}{Author Statement}
% ══════════════════════════════════════════════════════════════

This work was conducted independently by Eric Robert
Lawson under the OrganismCore research programme.
No institutional funding was received. No conflicts
of interest exist. No pharmaceutical company was
involved in the design of this proposal.
Tazemetostat is manufactured by Epizyme (now Ipsen).
Fulvestrant is manufactured by AstraZeneca (Faslodex)
and available in generic form. Neither company had
any involvement in this work.

This document is a scientific priority establishment
record and a trial proposal rationale. It is not a
registered clinical trial protocol. Actual trial
conduct requires ethics committee approval, regulatory
submission, institutional sponsorship, and full
protocol development by qualified clinical investigators.

% ══════════════════════════════════════════════════════════════
\section*{Contact}
\addcontentsline{toc}{section}{Contact}
% ══════════════════════════════════════════════════════════════

\begin{tabular}{ll}
Author        & Eric Robert Lawson \\
Organisation  & OrganismCore \\
Email         &
  \href{mailto:OrganismCore@proton.me}
       {OrganismCore@proton.me} \\
ORCID         &
  \href{https://orcid.org/0009-0002-0414-6544}
       {https://orcid.org/0009-0002-0414-6544} \\
Repository    &
  \href{https://github.com/Eric-Robert-Lawson/attractor-oncology}
       {github.com/Eric-Robert-Lawson/attractor-oncology} \\
Homepage      &
  \href{https://github.com/Eric-Robert-Lawson/OrganismCore}
       {github.com/Eric-Robert-Lawson/OrganismCore} \\
Date          & 2026-03-06 \\
\end{tabular}

\vspace{2em}

\begin{center}
\textit{``Two published papers confirm the mechanism.\\
No trial exists. Both drugs are approved.\\
170,000 patients per year.\\
The sequence is: tazemetostat first.
Then fulvestrant.\\
That is the trial.''}

\medskip
--- Eric Robert Lawson, March 6, 2026
\end{center}

% ══════════════════════════════════════════════════════════════
\end{document}
